\chapter{LM-IDNet Design and Implementation}
\label{ch:implementation}
This chapter describes the retained detector as an implemented procedure.
The DM family, Minka estimator, and LM numerical method are the inherited
foundations introduced in Chapter~\ref{ch:background}. The application adds
data contracts, normal-data calibration, a joint decision rule, guarded
learning, and a persistent runtime. The description follows the current
scoring implementation, including the count-first evaluation order; the
experimental records retain their individual source boundaries.

\section{System architecture}
\label{sec:architecture}

LM-IDNet separates preparation, fitting, calibration, inference, and
evaluation. Preparation translates mixed-device packet captures into one
device's ordered count windows. Fitting estimates the normal DM parameters.
Calibration establishes reference scales and cutoffs using a later benign
period. Inference consumes completed windows and produces events; optional
learning proposes changes for subsequent observations. Evaluation joins
saved events to attack annotations after inference has finished.

Figure~\ref{fig:pipeline} shows these dependencies. The separation between
fit and calibration data prevents the cutoff from being computed on the
same windows that selected the initial parameters. Temporal dependence and
development-set influence are handled separately by the experimental
protocol.

\begin{figure}[tbp]
\centering
\begin{tikzpicture}[
  box/.style={draw,rounded corners,align=center,text width=3.0cm,minimum height=1.2cm,font=\small},
  arrow/.style={-{Latex},thick},node distance=8mm and 6mm]
\node[box] (capture) {Source captures\\device selection};
\node[box,right=of capture] (window) {Validated counts\\UTC windows};
\node[box,right=of window] (fit) {Normal fit\\DM parameters};
\node[box,right=of fit] (cal) {Separate normal\\joint calibration};
\node[box,below=15mm of window] (input) {Later completed\\count windows};
\node[box,right=of input] (detector) {Shared detector\\fixed references};
\node[box,right=of detector] (events) {Events and\\update records};
\node[box,below=of detector] (state) {Validated candidate\\persistent pair};
\draw[arrow] (capture)--(window);
\draw[arrow] (window)--(fit);
\draw[arrow] (fit)--(cal);
\draw[arrow] (input)--(detector);
\draw[arrow] (cal.south)--++(0,-6mm)-|(detector.north);
\draw[arrow] (detector)--(events);
\draw[arrow] (detector)--(state);
\draw[arrow] (state.west)--++(-5mm,0)|-(detector.west);
\end{tikzpicture}
\caption{LM-IDNet data and state flow. The initial fit and calibration use
different normal captures. Runtime decisions are emitted before a window can
influence a candidate update. Labels join saved events only during offline
evaluation; they are absent from this runtime path.}
\label{fig:pipeline}
\end{figure}


One detector object owns the active parameters, matching calibration,
original reference, bounded buffers, and update counters. The batch runner
collects its per-window returns into finite experiment artifacts. The
streaming runner emits them as they become available. Both use the same
decision and update implementation. This prevents a separate deployment
implementation from drifting away from the model evaluated offline.

The command-line stages are explicit. A scoring request does not silently
retrain a model, and an evaluation request does not alter a threshold.
Configurations identify inputs, device semantics, model settings, output
locations, and adaptation policy. Configuration validation rejects unknown
or inconsistent settings before execution. Stored artifacts then provide
a second boundary: a valid configuration cannot make an incompatible
model and threshold into a meaningful pair.

Completed windows and JSON events define the supplied-window runtime. This
choice keeps the monitoring procedure inspectable; live collection,
transport durability, and response integration remain deployment work. The
detector's statistical
state is protected more strongly than ordinary research output files;
Section~\ref{sec:state} identifies that boundary.

\section{Capture validation, device filtering, and count windows}
\label{sec:preparation}

\subsection{Selecting the observed device}

The processor selects frames associated with the configured media access
control (MAC) address in either Ethernet direction. Matching ARP hardware
addresses are also considered. Filtering on this identity avoids relying
on a device's IP address remaining unchanged across captures. The device
name remains a human-readable identity in configurations and events; the
MAC address supplies the packet-selection rule. Source authentication
requires a separate mechanism.

The retained classification order is ARP, TCP, and then UDP. UDP packets
whose source or destination port is 1900 are assigned to SSDP; the
remaining UDP packets are assigned to ordinary UDP. Unsupported protocols
are excluded. Each retained packet increments one category, so the sum of
the four counts has a consistent meaning. This policy is versioned as
\path{device-mac-bidirectional-ssdp-exclusive-v1}. Changing it changes the
input semantics even if the resulting vector still has four entries.

Frame selection is performed before full packet decoding where possible.
This avoids spending decoding work on unrelated devices in large mixed
captures. Structural validation still inspects every source record. A fast
filter and a complete structural validator solve different preparation
problems.

\subsection{Window boundaries and observation states}

Let $t$ denote a packet's UTC timestamp in seconds from the Unix epoch and
let $\Delta$ be the trained duration in seconds. The containing window
starts at
\[
 b(t)=\Delta\left\lfloor\frac{t}{\Delta}\right\rfloor
\]
and represents $[b(t),b(t)+\Delta)$. A packet exactly at the right boundary
belongs to the next interval. The retained duration is $\Delta=600$~s.
Normalising timezone-aware timestamps to UTC before grouping provides one
time grid for preparation, annotations, and runtime validation.

Three observation states are distinguished. An observed window contains
counts of supported traffic; an observed-silent window contains a valid
zero vector; a missing window contains no count vector. Missing coverage is
excluded from fitting, calibration, and scoring. In the stream interface it
advances chronology without generating a decision, and gaps remain explicit.

This distinction also limits the security interpretation of silence.
The DM composition model assigns its conditional probability to the zero
vector when $N=0$, as explained in Chapter~\ref{ch:background}. The
retained upper count tails address high counts. An interruption of
communication therefore requires a separate availability signal, even when
the input correctly records an observed silent interval.

\subsection{Separate preparation gates}

Inventory records whether configured source files exist, their sizes and
checksums, and whether duplicated bytes need an explicit explanation.
Structural validation reads a classic-PCAP source to its physical end,
checking its global header, record boundaries, timestamp representation,
snapshot length, and captured/original length relationships. Recoverable
anomalies such as out-of-order timestamps are reported as warnings;
truncated records and impossible lengths fail with diagnostic evidence.
Structural validity comes from the complete record scan.

The classic-PCAP validator and the packet-reading path have different
format scopes. The dedicated structural validator covers classic PCAP;
the packet reader can also encounter PCAPNG through its decoder.
The UNSW exclusion of a truncated source is consequently a recorded
preparation decision, not evidence that all UNSW sources passed the same
classic-PCAP gate. The thesis retains this distinction when describing
data quality.

Dataset fingerprints hash source bytes and the policies that give those
bytes meaning: ordered categories, duration, timestamp/windowing rules,
and partition membership. Absolute clone locations and generation times
are excluded. A dataset identifier therefore changes when the experiment
definition changes, without changing merely because the repository moves
to another directory.

\section{Input contracts and temporal partitions}
\label{sec:contracts}

The saved model embeds its device identity, MAC address, category order,
duration, and counting policy. A vector can have the right dimension and
still be incompatible: swapping TCP and UDP entries, changing window
duration, or feeding another device's traffic changes its interpretation.
The runtime validates these declarations before assigning statistical
meaning to the counts. Compatibility is checked through the model and
calibration fingerprints as well as through their schemas.

At the trust boundary, counts must be non-negative integers. Boolean
values, numeric strings, fractions, and non-finite values are rejected.
The per-category counts and their total are bounded by $2^{63}-2$, leaving
room for adding one in factorial expressions. Numerical precision at the
upper end is enforced through explicit failures and recovery.

Time bounds must be ordered and match the trained duration. The streaming
boundary additionally enforces the UTC grid and rejects duplicates,
overlap, and backward movement. A valid input record declares both a
device and a source context; its source identifier need not encode a date.
The runtime validates the declaration; complete observation and
classification remain producer responsibilities.

Research partition access is also explicit. Fit and calibration stages
select only their declared captures. Ordinary scoring and evaluation use
the development role. A dedicated final runner selects the configured
final captures with additional checks on actual window chronology and
label coverage. Runtime operation uses the declared current source context,
independently of future captures and final-test identifiers.

Nominal capture dates and actual UTC times must not be confused. A file
named for a local calendar day can contain windows starting on the
preceding UTC date. Capture membership identifies the experimental role;
actual timestamps determine annotation alignment and overlap. The
separate temporal research branch supplies additional timeline audits beyond
the retained batch configuration checks. Its rolling-fold results remain a
separate evidence set from the selected detector.

\section{Model fitting and independent calibration}
\label{sec:calibration}

\subsection{Initial normal model}

For the $M$ observed fitting windows, initial category probabilities are
the pooled count proportions,
\[
 \widetilde p_k=\frac{\sum_{i=1}^{M}x_{ik}}
                       {\sum_{i=1}^{M}\sum_{j=1}^{K}x_{ij}},
 \qquad \alpha_k^{(0)}=c\,\widetilde p_k,
\]
where the configured initial concentration is $c=10$. Every category
must have positive aggregate count so that the initial alpha vector is
strictly positive. The implementation rejects an unusable initial matrix
rather than silently inventing an initial distribution.

Minka's update, Eq.~\eqref{eq:minka-update}, estimates the parameters
\cite{minka2000}. The row-wise kernel sum monitors convergence. The
reference settings stop when the absolute objective change is below
$2\times10^{-9}$ or the maximum of 100,000 iterations is reached. The
stored diagnostics distinguish convergence from merely reaching the
iteration limit. They retain initial/final objectives, parameter values,
iteration count, and elapsed fit time.

The stopping tolerance is an implementation criterion on a numerical
objective. Statistical uncertainty and convergence outside the available
iteration budget require separate analysis. The numerical comparison in
Chapter~\ref{ch:results} examines the behaviour of both backends under
that shared fitting procedure.

\subsection{Scaling three complementary features}

Let $g_1(\boldsymbol{x})=s_{\mathrm{DM}}(\boldsymbol{x})$ denote the
complete raw DM log-probability. Define
\[
 g_2(\boldsymbol{x})=-\log(1+N),\qquad
 g_3(\boldsymbol{x})=-\log(1+x_{\mathrm{ARP}}).
\]
The negative signs put unusually large counts into the same lower-tail
direction as unusually small composition scores. Both counts are already
available from the four-category input. No additional packet fields or
attack-labelled weights are required.

For each feature $j$, the fitting set supplies a median $c_j$ and an
interquartile range $d_j$. A zero fitting IQR is replaced by a unit scale.
The standardised feature and its joint summary are
\begin{equation}
 z_j(\boldsymbol{x})=\frac{g_j(\boldsymbol{x})-c_j}{d_j},
 \qquad J(\boldsymbol{x})=\min_{j\in\{1,2,3\}}z_j(\boldsymbol{x}).
\label{eq:joint-feature}
\end{equation}
Median/IQR scaling places heterogeneous features on reference scales;
calibration then estimates the shared empirical cutoff.
The unit-scale fallback is particularly relevant to a category that is
constant in the fitting data. Calibration still determines its effective
cutoff through the common joint rule.

\subsection{Joint calibration on separate normal windows}

If $\mathcal C$ denotes the benign calibration windows, let
\begin{equation}
 q_*=Q_{0.01}\{J(\boldsymbol{x}):\boldsymbol{x}\in\mathcal C\}.
\label{eq:joint-quantile}
\end{equation}
The reference calibration requires at least 30 observed windows and uses
linear interpolation for the sample quantile. Hyndman and Fan describe
the family of sample-quantile conventions to which this choice belongs
\cite{hyndman1996}. The fixed minimum is an engineering guard; tail
resolution still depends on the available calibration sample.

One joint quantile uses the observed dependence between the three
features, rather than selecting three thresholds independently and
assuming an overall false-alert probability. Its one-percent value is
an empirical operating point whose future false-positive rate depends on
the later traffic distribution.

For implementation, the composition cutoff is
$\tau=c_1+q_*d_1$, while the log-count cutoffs are
$u_N=-c_2-q_*d_2$ and $u_A=-c_3-q_*d_3$. The threshold artifact stores
these cutoffs, their reference scales, score convention, calibration
captures and time range, and a model fingerprint. Saving the cutoff
without this context would leave its numerical meaning ambiguous.

\section{Anomaly scoring and explanations}
\label{sec:joint-scoring}

\subsection{Signed margins and the decision}

For the active DM model, define the signed composition margin
$m_D=(s_{\mathrm{DM}}-\tau)/d_1$. The count margins are
\[
 m_N=\frac{u_N-\log(1+N)}{d_2},\qquad
 m_A=\frac{u_A-\log(1+x_{\mathrm{ARP}})}{d_3}.
\]
With all three features evaluated, the static combined score is
\begin{equation}
 S(\boldsymbol{x})=\min\{m_D,m_N,m_A\},
 \qquad \widehat y=\mathbf{1}\{S(\boldsymbol{x})<0\}.
\label{eq:combined-decision}
\end{equation}
This is equivalent to comparing Eq.~\eqref{eq:joint-feature} with
Eq.~\eqref{eq:joint-quantile}. The comparison is strictly below zero, so
equality at a cutoff remains normal. The minimum makes any evaluated
negative margin sufficient evidence for an alert, without interpreting
that margin as an attack probability.

In adaptive operation, the original DM and its original calibrated
scale/cutoff remain detection references. When composition is evaluated,
its margin is the minimum of the active and original DM margins.
The original count cutoffs also remain fixed. Consequently, a candidate
cannot remove a decision supported by an original reference merely by
making the active model more tolerant. This can preserve desirable
evidence, but can also preserve an original false alert after benign
change. Reference refresh requires an explicit new model and calibration.

\subsection{Count-first evaluation and numerical recovery}

The current static and conservative paths evaluate the inexpensive count
margins first. A negative count margin already determines the alert in
Eq.~\eqref{eq:combined-decision}. The implementation can therefore emit
that decision without evaluating composition or learning from the window.
Its event contains the evaluated count margins and a null composition
score. Otherwise the DM is evaluated and the composition margin joins
the minimum. The historical adaptive policy retains complete DM scoring
because its recent-score buffer depends on those values.

This evaluation order preserves the binary rule but need not preserve the
numerical minimum that would have been obtained by evaluating every
feature. A skipped DM margin could have been more negative. Event severity
and score rankings can consequently change while anomaly decisions remain
equal. The later scoring-order regression is reported separately from
the original full-score final replay for this reason.

If the native precision guard rejects a required DM calculation, the
implementation recomputes the same complete probability with the installed
high-precision library and logs the recovery. The model, coefficient
convention, and cutoff remain fixed. Typed errors and recovery delimit the
requested native precision and accepted integer range.

\subsection{What an event communicates}

An event identifies the device, capture context, UTC bounds, counts,
score convention, evaluated margins, decision cutoff, and model
fingerprint. It also records the expected proportions and residuals
\[
 r_k=\begin{cases}x_k/N-p_k,&N>0,\\-p_k,&N=0.\end{cases}
\]
The silent-window convention uses zero observed proportions. It is a
consistent diagnostic convention for the empty sample.

For a combined event, severity is $\max(0,-S)$ using the score actually
evaluated. It expresses distance beyond a cutoff on the recorded reference
scale. A large residual points to an over- or under-represented category;
attack attribution requires additional evidence. A count-first event
reports only the evidence actually evaluated. The event contract makes
that absence visible to a reviewer.

\subsection{A worked probability and decision example}
\label{sec:worked-score}

A small synthetic example makes the scoring convention concrete by
illustrating Eq.~\eqref{eq:dm-pmf}. Consider four categories
with $\boldsymbol{\alpha}=(1,1,1,1)$ and a window
$\boldsymbol{x}=(2,1,0,0)$, so $N=3$ and $\alpha_0=4$.
The complete count-vector probability is
\begin{align*}
 P(\boldsymbol{x}\mid3,\boldsymbol{\alpha})
 &=\frac{3!}{2!1!0!0!}
   \frac{\Gamma(4)}{\Gamma(7)}
   \frac{\Gamma(3)\Gamma(2)\Gamma(1)\Gamma(1)}{\Gamma(1)^4}\\
 &=3\cdot\frac{6}{720}\cdot2=\frac{1}{20}.
\end{align*}
Its complete log probability is approximately $-2.995732$.
Omitting the multinomial coefficient would instead give $1/60$,
or approximately $-4.094345$ in log units. The discrepancy comes from
calculating a different probability event.

Under this particular uniform Dirichlet prior, all twenty possible
four-category count vectors summing to three have equal probability.
For example, $(3,0,0,0)$ has coefficient one and a different Gamma
product, but the same complete probability $1/20$. A rule omitting
the coefficient would make these vectors appear to have different
probabilities. This shows why comparing implementations requires an
explicit count-vector convention, even though the coefficient can
be omitted when maximising parameters for fixed data.

For an illustrative calibration, suppose the composition cutoff is
$\tau=-3.2$ with $d_1=0.4$. Let the total log-count cutoff be
$u_N=\log11$ with $d_2=0.5$, and the ARP log-count cutoff be
$u_A=\log6$ with $d_3=0.5$. These are chosen numbers, not the saved
project thresholds. Since the total is three and the ARP count zero,
the margins are
\begin{align*}
 m_D&=\frac{-2.995732-(-3.2)}{0.4}\simeq0.511,\\
 m_N&=\frac{\log11-\log4}{0.5}\simeq2.023,\\
 m_A&=\frac{\log6-\log1}{0.5}\simeq3.584.
\end{align*}
All margins are positive, so their minimum is positive and the
window remains normal. A zero cutoff in margin space is compatible
with a negative cutoff in log-probability space. Confusing the two
thresholds would produce a different decision rule.

The example also separates decision and learning. The observation
is non-alerting, but its composition margin is below the conservative
admission margin of one. It is therefore excluded from that learning
buffer even if its volume passes the additional admission limit. Adaptation
eligibility uses the stricter admission rule.

Count-first execution preserves the logical decision by stopping
as soon as a negative count margin is known. It need not compute a
complete composition explanation for that event. The retained
examples in Section~\ref{sec:event-cases} include a real window with
both composition and count violations, for which this distinction
changes the reported score while leaving the alert unchanged.

\section{Native LM backend}
\label{sec:native-integration}

The supplied numerical code is retained separately from the application's
production native copy. The integration preserves the LM approximation
and its attribution to Languasco and Migliardi \cite{languasco2023}. Its
changes concern validation, matrix semantics, coefficient loading, memory
lifecycle, zero-count increments, and the Python/native boundary.
The thesis claims these integration decisions, not the mathematical
invention of the LM method.

Python sends a contiguous matrix and alpha vector to the compiled
library. A native matrix entry point evaluates each row's kernel and
returns their sum. Numerical setup that depends on alpha, category count,
and requested precision can be reused across rows within that call.
Observed counts still enter each row independently. Replacing the matrix
with its summed count vector would change the objective, as established
in Chapter~\ref{ch:background}.

The earlier adapter repeated the shared setup for each row. Moving that
work outside the row loop removes an integration cost while retaining
row-wise semantics. The corresponding timing comparison and its historical
boundary are reported in Chapter~\ref{ch:results}; performance conclusions
remain attached to the measured workloads.

The adapter validates library availability, dimensions, finite positive
parameters, and the returned result. Coefficient-file reads and allocations
are checked, and workspaces are released on successful and failed calls.
Native failures are translated to application errors rather than
terminating the interpreter. Calls are serialised around global state
inherited from the supplied implementation. This protects the existing
kernel contract without claiming it is an independently re-entrant native
library.

The SciPy implementation supplies a reference kernel through
\texttt{gammaln}, while SciPy's digamma remains in the common estimator
\cite{virtanen2020}. The count-dependent coefficient is added outside
either kernel. Numerical validation therefore covers both backend agreement
and the complete score seen by calibration. The offline image contains
the compiled library and coefficient files, enabling runtime start from
bundled resources.

\section{Legacy and conservative adaptation}
\label{sec:adaptation}

\subsection{Mode, policy, and update ordering}

A mode determines what may change: static operation changes neither
parameters nor threshold, threshold adaptation changes its cutoff and
scale, and periodic refitting proposes a complete parameter/calibration
pair. A policy determines which observations and candidates are permitted
to cause those changes. Both the historical legacy policy and the
conservative policy use the same per-window processing path.

The current observation is scored before any admission or candidate
activation. A promoted update can affect later windows only. Missing
observations do not enter learning buffers or the scored-window counter.
The ordinary attempt schedule is expressed in observed windows, not
elapsed wall time. During a long gap, it does not silently advance as if
traffic had been observed.

The reference ten-minute configuration uses a bounded buffer of 288
observations and an attempt cadence of 144 observations. Periodic
refitting normally requires 144 accepted observations. These correspond
to 48, 24, and 24 observed hours at ten minutes, respectively; a sparse
accepted buffer can span a much longer chronological period. The historical
288/240 experiment changes cadence and minimum acceptance while retaining
the same buffer capacity. Results identify this schedule separately.

\subsection{Admission to a conservative learning buffer}

Let $s_t$ and $s_0$ denote raw DM scores under the active and original
models. Write $\tau_t,I_t$ and $\tau_0,I_0$ for their cutoffs and positive
reference scales. With confidence multiplier $h=1$, conservative admission
requires
\begin{equation}
 \begin{split}
  s_t(\boldsymbol{x})&\geq\tau_t+hI_t,\\
  s_0(\boldsymbol{x})&\geq\tau_0+hI_0,\\
  N&\leq V,\qquad m_N\geq h,\qquad m_A\geq h,
 \end{split}
\label{eq:admission}
\end{equation}
where $V$ is the 0.99 empirical quantile of benign calibration totals.
The volume cap $V$ restricts learning and is distinct from the joint
detection cutoff $u_N$. The cap can reject a window that is not an alert.
Both active and original composition evidence must clear a confidence
margin, and count margins must be safely positive.

Equation~\eqref{eq:admission} admits observations using unlabelled
statistical evidence; benign status remains unknown. Attacks that resemble
the observed normal features can still pass. The poisoning and contamination
literature motivates explicit admission controls
\cite{kloft2012,cretu2008}; the specific DM rule is evaluated empirically.
Admission counters distinguish excessive volume from insufficient
confidence so that the restriction's empirical effect remains inspectable.

\subsection{Candidate construction and validation}

Threshold-only candidates score accepted counts under the active DM and
recalculate the lower quantile and IQR. Refitting uses the same estimator
as initial training, applied to accepted counts with one synthetic
single-count row per category. The four smoothing rows ensure each
category has positive total count for initialisation. They are a fitting
device excluded from observed-window and detection-label counts.

The candidate must be finite, positive, dimensionally compatible, and
converged. Its kernel objective on the refitting matrix must not be worse
than that of the active model beyond numerical tolerance. Conservative
validation additionally imposes
\begin{equation}
 \frac{1}{1.25}\leq\frac{\alpha'_0}{\alpha_{0,t}}\leq1.25,
 \quad \max_k|p'_k-p_{k,t}|\leq0.05,
 \quad |\tau'-\tau_t|\leq2I_0.
\label{eq:movement-bounds}
\end{equation}
The candidate scale must be finite and positive. Threshold and model are
validated together, and failure rejects the entire conservative pair.
There is no fallback that applies an otherwise unrelated threshold change
after a conservative refit is rejected.

These tests control particular movements. The likelihood check uses the
fitting buffer rather than an independent candidate holdout. The
concentration and proportion bounds compare successive active models,
so a sequence of individually bounded moves need not have a fixed
cumulative bound. The threshold bound uses the original scale. Poisoning
resistance and future false-alert stability therefore remain empirical
outcomes.

\subsection{The preserved historical policy}

Legacy threshold adaptation admits all recent scored DM values, including
anomalous ones. Legacy refitting uses an original-baseline confidence
gate, convergence, and likelihood checks, but lacks the newer volume and
movement restrictions. A rejected legacy refit may still fall back to a
threshold update. With the retained combined detector, fixed detection
references remain in effect, while historical DM-only artifacts preserve
their earlier scoring contract.

These distinctions permit a fair description of negative results.
Learning from all recent scores can move a cutoff toward attack traffic;
a converged refit can fit contaminated observations well. Conversely,
restrictive guards can preserve starting decisions by rejecting every
candidate. A rejection verifies enforcement of a configured condition,
not successful adaptation. Controlled promotion tests establish that the
software can activate an allowed candidate; empirical traffic experiments
must establish whether doing so is useful.

\begin{figure}[tbp]
\centering
\fbox{\begin{minipage}{0.92\textwidth}\small
\textbf{One completed-window operation}\par
\begin{enumerate}
 \item Validate identity, categories, duration, values, and observation state.
       Return no event for missing coverage.
 \item Compute count margins. In static/conservative operation, emit an
       alert immediately if a count margin is negative.
 \item Otherwise evaluate the required DM scores, combine active/original
       evidence, and construct the event for the current window.
 \item In adaptive operation, test admission and update bounded buffers.
 \item At an eligible attempt boundary, construct a threshold or refit
       candidate using the accepted data.
 \item Check sample support, numerical validity, convergence, likelihood,
       and configured movement limits. Retain the active pair on rejection.
 \item Persist a valid complete pair before activating it. Return the
       current event and the candidate's accepted/rejected audit outcome.
\end{enumerate}
\end{minipage}}
\caption{Processing order shared by batch and streaming operation. The legacy
policy retains full DM scoring for its historical recent-score buffer. An
activated candidate affects subsequent observations, not the decision already
constructed for the current window.}
\label{fig:window-procedure}
\end{figure}

An early count violation returns from the operation after recording
admission rejection in an adaptive mode, before learning-buffer insertion
or fitting, even if that scored-window index is an
ordinary scheduled boundary. Later eligible boundaries still use
the observed-window counter. The regression evidence establishes
equal decisions and final states on its tested histories; it does
not claim that every possible future attempt history is identical
to the older full-score execution order.

\section{Persistence, restart, and rollback}
\label{sec:state}

Persisting alpha alone would leave the detector vulnerable to combining a
new model with an old cutoff. A snapshot therefore stores a complete
active pair and its original references, input contract, calibration and
adaptation settings, estimator settings, and the learning-volume limit.
The original threshold retains fixed count cutoffs. Restoring a DM-only
snapshot preserves its DM-only scoring contract.

Versions are immutable and identified separately from parameter
fingerprints. A pointer file names current and previous versions, and a
bounded history records recent decisions. A nonblocking Linux file lock
allows only one state writer or rollback operation at a time. Locking
protects this local state directory. Distributed replicas would require a
different coordination protocol.

Figure~\ref{fig:state-lifecycle} identifies the activation boundary. The
writer creates a temporary file beside its destination, flushes and
synchronises it, and reloads the snapshot through validation. The immutable
version is published before the active pointer is atomically replaced.
Failure before pointer replacement leaves the former pair active.
Failure to confirm directory durability after replacement is a committed
state error: the runtime stops instead of describing the update as rejected.
Restart then validates the state actually referenced by the pointer.

\begin{figure}[tbp]
\centering
\begin{tikzpicture}[node distance=7mm,
 box/.style={draw,rounded corners,align=center,text width=10.3cm,minimum height=8mm,font=\small},
 arrow/.style={-{Latex},thick}]
\node[box] (candidate) {Candidate parameters, threshold, scales, and original references};
\node[box,below=of candidate] (validate) {Validate the complete candidate; write, flush, and reload its snapshot};
\node[box,below=of validate] (version) {Publish an immutable version file; keep the current pair active};
\node[box,below=of version,very thick] (pointer) {Atomically replace the current-version pointer: activation boundary};
\node[box,below=of pointer] (restore) {Synchronise directory; activate the same pair in memory};
\node[box,below=of restore] (retain) {Retain current and previous versions; bound diagnostic history};
\foreach \a/\b in {candidate/validate,validate/version,version/pointer,pointer/restore,restore/retain}
 \draw[arrow] (\a)--(\b);
\end{tikzpicture}
\caption{Accepted-state publication. Failure before pointer replacement preserves
the former active pair. Uncertain durability after replacement stops the writer;
restart validates the committed reference rather than reporting a rejected update.}
\label{fig:state-lifecycle}
\end{figure}


Cleanup keeps active and previous snapshots and the most recent 100
history decisions. Managed obsolete versions and abandoned temporary
files may be removed only after reference validation. Unrelated files are
left alone. Cleanup or logging failure is distinct from reversing a
successful activation. The guarantees concern accepted detector state;
research reports and event arrays use their own finite-run persistence.

Restoration validates device semantics and detector settings. An existing
corrupt state fails instead of resetting silently to the initial model.
Prepared edge state embeds the required calibration evidence, allowing
startup without research captures. Buffers, scheduling counters, and input
positions restart empty. The external producer therefore owns input position
and duplicate-delivery handling.

Rollback requires a previous accepted version and exclusive writer access.
It restores that complete pair rather than mixing selected fields from
two versions. An operator must stop the current writer first. This
mechanism supports recovery from a known undesirable accepted change;
the operator supplies the reason for rollback.

\section{Bounded streaming runtime and Docker simulation}
\label{sec:streaming}

The stream interface reads one UTF-8 JSON record per physical line. Each
record declares the device, source context, and a nested completed window.
The 65,536-byte maximum includes the line ending and is enforced while
reading. Malformed JSON, non-finite numbers, missing fields, incompatible
semantics, and chronology violations terminate with a line-numbered error.
A valid final record is accepted with or without a terminating newline.

After processing, an event is serialised and flushed before the next
input read. Input and output are synchronous: a slow consumer pauses
further reads rather than causing an unbounded queue. The detector keeps
bounded deques, active/original state, and counters. Batch result arrays
grow with the finite research experiment; this deliberate collection is
not part of the streaming memory bound.

Signals distinguish idle reading from an active detector operation.
SIGINT or SIGTERM stops an idle reader; during scoring or persistence,
shutdown is deferred until the current operation completes. This avoids
interrupting activation mid-transaction. The service's stop grace period is
30~s before Docker's final termination mechanism.

The packaged simulation uses a read-only Linux/amd64 image, a one-CPU
quota, 512~MiB of memory, a 16-MiB temporary filesystem, and a 64-process
limit. Configuration is read-only; prepared state and bounded logs are
writable using the host's numeric identity. Networking is disabled, and
the image starts without compiling or downloading dependencies.
Application logs retain a current 5-MiB file and two backups. The local
Docker log driver retains at most three 5-MB files, in addition to its
own metadata and compression behaviour.

Normal end-of-file exits successfully and leaves the finite service
stopped. Failures can trigger up to three retries; an explicit stop stays
stopped. A collector-owned FIFO can supply records, but that collector
owns reconnection, input position, live capture, and durable event
acknowledgement. An update can be
committed before its associated event reaches a failed consumer. These
limits are compatible with the supplied-window simulation and are retained
when interpreting its resource measurements.

\section{Computational and storage boundaries}
\label{sec:computational-boundaries}

The compact feature dimension reduces data volume. Preparation still reads the source
capture and validates its records, even when most frames belong to
other devices. Cost depends on captured traffic and decoding, not
merely on the four numbers eventually stored per window. Offline
fitting retains an $n\times K$ matrix and repeatedly evaluates
digamma terms and likelihood checks until convergence or the
iteration limit. Iteration count and special-function evaluation
cost remain input-dependent.

Inference retains active and original parameter/calibration objects
of size proportional to $K$, bounded learning buffers containing at
most $B$ count vectors, and bounded diagnostic history of length
at most $H$. Their dominant stored numeric content is proportional
to $BK+H$, rather than the total number of stream records. Here
$K=4$, $B=288$ for the reference policy, and persisted history is
limited to 100 records. The reader has its per-line bound. This
describes structures controlled by the application; measured RSS also
includes the interpreter and native allocator.

An ordinary window evaluates count features and, where needed,
composition references. A scheduled refit constructs and fits a
bounded candidate matrix, including the pseudo-observations. Its
time can dominate that window even though its storage remains
bounded and the candidate is rejected. Admission can delay an
attempt, while convergence checks can make an actual attempt
expensive. The resource study therefore reports both mean and worst-case
operation time.

Streaming serialises one result and waits for its consumer before
reading another record, keeping event memory bounded.
Batch evaluation deliberately saves all finite experiment events
for later label joining. Disk space for captures, experiment
outputs, and external event retention is separate from bounded
learning state. These boundaries explain why the resource study
measures both peak memory and refit latency. Raw packet collection requires
a separate end-to-end measurement.
